\chapter[2001 --- Hartwell et al.]{2001: Leland H. Hartwell, Timothy Hunt, Sir Paul M. Nurse}
\label{ch:2001_Hartwell}

\section*{Main Contribution}

\textit{Discovered key regulators of the cell cycle (cyclins and CDKs), explaining how cells control their division---critical for understanding cancer.}

\section{Official Citation}

for their discoveries of key regulators of the cell cycle

\section{Background and Historical Context}

All living organisms are built from cells, and all cells arise by division. The cell cycle---the ordered sequence of events by which a cell replicates its DNA and divides into two daughter cells---is one of the most fundamental processes in biology. Dysregulation of the cell cycle is a hallmark of cancer, and understanding the molecular mechanisms that control cell division has been one of the major goals of biomedical research in the twentieth century. The discovery of key regulators of the cell cycle by Leland H. Hartwell, Timothy Hunt, and Sir Paul M. Nurse, recognized with the 2001 Nobel Prize in Physiology or Medicine, provided the molecular framework for understanding how cells grow, replicate their DNA, and divide, and has had transformative implications for cancer biology and therapy.

Leland Harrison Hartwell was born on October 7, 1939, in Los Angeles, California. He studied physics at the California Institute of Technology (Caltech) before switching to biology, earning his Ph.D.\ from the Massachusetts Institute of Technology (MIT) in 1964 under the supervision of Terrell Edwards, working on the genetics of the yeast Saccharomyces cerevisiae. Hartwell joined the faculty at the University of Washington in Seattle in 1965, where he conducted his Nobel Prize--winning research.

Timothy (Tim) Hunt was born on February 19, 1943, in Neston, Cheshire, England. He studied natural sciences at the University of Cambridge and earned his Ph.D.\ in 1968 under the supervision of Frederick Sanger, working on the synthesis of hemoglobin. Hunt then spent a year at the Institute of Molecular Biology in Cambridge, where he began studying protein synthesis during the cell cycle. He joined the staff of the Medical Research Council (MRC) Laboratory of Molecular Biology in Cambridge in 1976, where he conducted his Nobel Prize--winning work.

Sir Paul Maxime Nurse was born on January 25, 1949, in Norfolk, England. He studied biochemistry at the University of Birmingham and earned his Ph.D.\ from the University of East Anglia in 1973, working on the genetics of the yeast Schizosaccharomyces pombe. Nurse joined the MRC Cell Biology Unit in London and later became director of the Imperial Cancer Research Fund (now Cancer Research UK) in London. He served as president of Rockefeller University from 2003 to 2011 and became director of the Francis Crick Institute in London in 2011.

The intellectual context of their work was the recognition that the cell cycle is a precisely regulated process governed by a series of molecular checkpoints. The two principal events of the cell cycle---DNA replication (S phase) and mitosis (M phase), separated by gap phases G1 and G2---must occur in the correct order and with precise timing. The mechanisms that ensure this order and timing were unknown until Hartwell, Hunt, and Nurse made their discoveries using yeast genetics and biochemistry.

\section{The Discovery/Contribution}

\subsection{Leland Hartwell and Cell Division Cycle (cdc) Genes}

Hartwell used the budding yeast S. cerevisiae as his experimental model, exploiting the power of yeast genetics to identify genes required for cell division. He isolated temperature-sensitive mutants---strains that grow normally at a low (permissive) temperature but fail to divide at a high (restrictive) temperature---and screened them for defects in specific stages of the cell cycle. Through systematic genetic analysis, Hartwell identified more than 30 ``cell division cycle'' (cdc) genes, each encoding a protein essential for progression through a specific step in the cell cycle.

Two of Hartwell's discoveries were particularly important. First, he identified the concept of ``cell cycle checkpoints''---control mechanisms that ensure the completion of one cell cycle event before the next is initiated. For example, the DNA damage checkpoint pauses the cell cycle when DNA is damaged, allowing time for repair before replication proceeds. The spindle checkpoint ensures that all chromosomes are properly attached to the mitotic spindle before anaphase begins. These checkpoint mechanisms are fundamental to the fidelity of cell division and are often disrupted in cancer.

Second, Hartwell identified the CDC28 gene, which encodes a protein kinase that controls the transition from G1 to S phase---the ``start'' point of the cell cycle. The CDC28 protein was the first cell cycle regulator to be identified by genetic analysis and was later shown to be the prototype of the cyclin-dependent kinase (CDK) family.

\subsection{Timothy Hunt and Cyclins}

Hunt's contribution was the discovery of cyclins---proteins whose concentration rises and falls in a regular pattern during the cell cycle. In 1982, while studying protein synthesis during the early embryonic cell cycles of the sea urchin Arbacia punctulata, Hunt observed that a small number of proteins are synthesized and then destroyed in a periodic fashion, with peaks of synthesis corresponding to specific phases of the cell cycle. One of these proteins, which he named ``cyclin,'' accumulated during interphase and was rapidly degraded at the onset of mitosis.

Hunt showed that cyclin degradation is essential for the exit from mitosis: if cyclin is not destroyed, cells cannot complete mitosis and divide. He further showed that cyclin acts by binding to and activating CDKs, forming cyclin--CDK complexes that phosphorylate target proteins and drive progression through the cell cycle. Different cyclins, produced and degraded at different times, activate different CDKs at different stages of the cell cycle, creating the temporal program that governs cell division.

Hunt's discovery was made using sea urchin embryos, which undergo rapid, synchronous cell cycles in the absence of cell growth, providing a simplified system for studying cell cycle oscillations. The identification of cyclins in sea urchins was soon extended to other organisms---from yeast to humans---revealing the universal conservation of the cyclin--CDK regulatory mechanism.

\subsection{Paul Nurse and CDKs}

Nurse's contribution was the identification and characterization of cyclin-dependent kinases (CDKs) in the fission yeast S. pombe. Working independently of Hartwell, Nurse isolated temperature-sensitive mutants of S. pombe that failed to divide, and identified the CDC2 gene, which encodes a protein kinase essential for the G1/S and G2/M transitions. Nurse showed that the CDC2 protein is a kinase that is activated by cyclins and that phosphorylates target proteins to drive cell cycle progression.

Nurse's primary contribution was the demonstration that the CDC2 kinase from fission yeast is functionally equivalent to the CDC28 kinase from budding yeast and to the CDK1 kinase from human cells. This cross-species complementation experiment proved that the molecular mechanisms controlling the cell cycle are conserved across all eukaryotic organisms---from yeast to humans. It also demonstrated that CDKs are the master regulators of the cell cycle, acting as the catalytic subunits that are activated by cyclins and that phosphorylate the substrates that drive cell cycle transitions.

Nurse also identified the concept of the ``mitotic trigger''---the abrupt activation of CDK1--cyclin B at the G2/M transition, which triggers the dramatic events of mitosis: chromosome condensation, nuclear envelope breakdown, spindle formation, and sister chromatid separation. He showed that the mitotic trigger operates as a switch-like, all-or-none mechanism, ensuring that mitosis is initiated rapidly and completely once the cell has passed the checkpoint.

\section{Impact on Modern Medicine}

The discoveries of Hartwell, Hunt, and Nurse have had a transformative impact on cancer biology and therapy.

\subsection{Cancer Biology}

Cancer is fundamentally a disease of uncontrolled cell division. The cell cycle regulators discovered by Hartwell, Hunt, and Nurse are frequently mutated or dysregulated in cancer. The tumor suppressor gene RB (retinoblastoma protein) controls the G1/S checkpoint by inhibiting CDK--cyclin complexes. The tumor suppressor gene p53 activates the DNA damage checkpoint by inducing the expression of CDK inhibitors such as p21. The oncogene cyclin D1 is amplified or overexpressed in many cancers, driving inappropriate cell cycle progression.

Understanding the molecular mechanisms of cell cycle control has enabled the classification of cancers based on their molecular profiles and has identified new targets for therapeutic intervention. CDK inhibitors---small molecules that block the activity of specific cyclin--CDK complexes---are now in clinical use. Palbociclib, ribociclib, and abemaciclib, which inhibit CDK4/6, are approved for the treatment of hormone receptor-positive breast cancer and represent a new class of targeted anticancer agents.

\subsection{Cancer Therapy}

The cell cycle checkpoints discovered by Hartwell are also exploited by existing cancer therapies. Many chemotherapeutic drugs (including cisplatin, doxorubicin, and etoposide) work by inducing DNA damage, which activates the DNA damage checkpoint and triggers cell cycle arrest and apoptosis. The efficacy of these drugs depends on the integrity of the checkpoint mechanisms. Understanding how checkpoints are disrupted in cancer has led to the development of combination therapies that target both the cell cycle machinery and the checkpoint pathways.

\subsection{Stem Cell Biology and Tissue Regeneration}

The cell cycle regulators also play essential roles in stem cell biology. The balance between self-renewal and differentiation in stem cells is controlled by cyclin--CDK activity, and perturbation of this balance can lead to either stem cell exhaustion or tumor formation. Understanding these mechanisms is essential for the development of stem cell-based therapies for degenerative diseases.

\section{Legacy}

Leland Hartwell became president and director of the Fred Hutchinson Cancer Research Center in Seattle and later became director of the Center for Sustainable Health at Arizona State University. He received numerous honors, including the Albert Lasker Basic Medical Research Award in 1998. Hartwell has been a passionate advocate for the prevention of chronic diseases through lifestyle modification.

Timothy Hunt continued his work at the MRC Laboratory of Molecular Biology in Cambridge and later at the Francis Crick Institute. He received the Albert Lasker Basic Medical Research Award in 1998 and the Royal Medal in 2006. Hunt continues to study the cell cycle and has been a mentor to numerous scientists who have gone on to make major contributions to the field.

Paul Nurse served as president of Rockefeller University and then became founding director of the Francis Crick Institute in London. He received a knighthood in 1999 and numerous other honors, including the Albert Lasker Basic Medical Research Award in 1998. Nurse continues to study the cell cycle and has been a vocal advocate for science education and public engagement with science.

Together, Hartwell, Hunt, and Nurse revealed the molecular engine that drives cell division. Their discoveries of cyclins and cyclin-dependent kinases have provided the framework for understanding how cells grow, divide, and are controlled in health and disease. The ongoing application of their insights to cancer biology, stem cell research, and regenerative medicine ensures that their legacy will continue to grow.


\section{Modern Developments}

Hartwell, Hunt, and Nurse's cell cycle work has led to modern cancer therapy. Understanding of cyclins and CDKs has led to CDK4/6 inhibitors (palbociclib, ribociclib) for breast cancer. Understanding of checkpoint kinases has led to ATR inhibitors for DNA damage repair-deficient cancers. Understanding of cell cycle control has informed development of CDK inhibitors for other cancers. Understanding of mitotic machinery has led to taxanes and vinca alkaloids for cancer chemotherapy.


\section{Bibliography}

\begin{thebibliography}{9}
\bibitem{nobel-2001}
Nobel Prize Outreach, ``The Nobel Prize in Physiology or Medicine 2001,''\emph{NobelPrize.org}.\\
\url{https://www.nobelprize.org/prizes/medicine/2001/summary/}

\bibitem{lecture-2001}
Leland H. Hartwell, ``Nobel Lecture,''\emph{NobelPrize.org}. Winner-authored primary source; downloadable from the official Nobel Prize archive.\\
\url{https://www.nobelprize.org/prizes/medicine/2001/hartwell/lecture/}
\end{thebibliography}
